\chapter{Conclusiones}

Este trabajo validó el diseño estructural de la pala de un aerogenerador offshore flotante de eje vertical mediante simulación aerodinámica en QBlade, conforme a los criterios de la norma IEC 61400-1~\parencite{IEC61400-1}. Las conclusiones se organizan en torno al cumplimiento de los objetivos planteados, las limitaciones del alcance y las líneas de trabajo que se abren a partir de los resultados obtenidos.

\section{Cumplimiento de los objetivos}

Respecto del primer objetivo específico, el comportamiento aerodinámico de la pala quedó caracterizado en la totalidad del rango operacional, y las 102 simulaciones LLFVW entregaron las series temporales de cargas que el objetivo exigía. Lo relevante para lo que sigue es que esas cargas llegan respaldadas, ya que cada eslabón de la cadena que las produce quedó contrastado con una referencia ajena al propio cálculo: las polares del NACA 0018 contra dos conjuntos de datos experimentales, el TSR óptimo del barrido DMST contra la solidez del rotor como parámetro puramente geométrico, y la curva de potencia resultante contra las turbinas H-rotor documentadas de escala comparable.

De ese objetivo se desprende también la velocidad máxima de rotación, que el fabricante no especifica y que se estimó en $42$~rpm. El valor resiste el contraste con las máquinas de referencia (Tabla~\ref{tab:comparacion_adim}), pero su consecuencia estructural no se hereda de ellas, ya que el menor radio del rotor en estudio eleva la aceleración centrífuga un $29\%$ sobre la de la T1-turbine en su propio límite de diseño. Siendo esta la solicitación que domina el estado tensional, conviene subrayar que se trata de un parámetro estimado, cuya confirmación queda pendiente de que la empresa libere el valor de diseño efectivo.

En cuanto al segundo objetivo, se desarrolló un modelo de elementos finitos de la pala con condiciones de borde y estados de carga equivalentes al escenario operacional, esto es, empotramiento en las conexiones con los struts, aceleración centrífuga de $20{,}4\,g$ mayorada por el factor parcial del caso, y fuerzas aerodinámicas de diseño aplicadas por estación conservando el signo del instante del que provienen. Dos análisis de convergencia de malla, uno sobre la tensión y otro sobre el autovalor de pandeo, acreditan que los resultados no dependen de la discretización.

De ese modelo se desprende una precaución para el diseño. La tensión máxima aparece en la zona de conexión con los brazos, que es también la que hubo que suponer ante la falta de un plano de la unión, de modo que el factor de reserva obtenido es sensible al detalle constructivo definitivo. Con un margen del $14\%$, ese detalle debe verificarse cuando se disponga de su geometría real.

El tercer objetivo se cumplió al verificar la integridad estructural bajo los tres casos de carga de producción de energía de la norma. La resistencia última se satisface con un factor de reserva de $1{,}26$, la deflexión no supera el $1{,}3\%$ de la luz entre conexiones y el daño por fatiga en 20 años, $D_{20} = 6{,}00 \times 10^{-3}$, queda por debajo del límite unitario con más de dos órdenes de magnitud de margen. La estabilidad, en cambio, no se cumple, y no por un margen estrecho, ya que el multiplicador de carga del primer modo vale $0{,}569$ frente al umbral de $1{,}20$ que la norma fija para el pandeo global de cascarones curvos. Falta un factor $2{,}11$, y al no alcanzar el multiplicador siquiera la unidad, el modelo predice inestabilidad bajo la propia carga de diseño. Es, con claridad, el criterio que gobierna el diseño.

Ese margen a fatiga merece una precisión. Por la planitud de la curva S-N del laminado, el daño escala con la amplitud de tensión elevada a $22{,}2$, de modo que los dos órdenes de magnitud de holgura sobre $D_{20}$ equivalen a un margen de apenas un $20\%$ sobre la tensión. La fatiga no es dimensionante, pero tampoco dispone del colchón que la cifra sugiere a primera vista, y su verificación queda condicionada a la precisión con que se calibre la relación entre momento y tensión en la sección crítica.

Que el criterio incumplido sea la estabilidad y no la resistencia última no es lo que cabría anticipar, y su explicación se desarrolla en la Sección~\ref{sec:resultados_fem}. De ella se desprende una conclusión con consecuencias para el diseño, y es que el margen de pandeo obtenido debe leerse como una señal de que la sección requiere revisión, y no como una cuantificación exacta del déficit, porque los supuestos del modelo sobre la distribución del espesor y sobre la orientación del refuerzo recaen precisamente sobre la magnitud que ese criterio evalúa.

Con los tres objetivos específicos cumplidos, el objetivo general queda alcanzado en su propósito de someter el diseño a la validación normativa. El resultado de esa validación es matizado, ya que la pala es válida en resistencia última, fatiga y deflexión bajo condiciones de funcionamiento nominal según IEC 61400-1, pero la verificación de estabilidad muestra que la configuración analizada, piel de $8$~mm con larguero interior del mismo espesor, queda por debajo del margen normativo de pandeo, por lo que la validez del diseño queda condicionada a la rigidización de los paneles de piel del vano central.

Del análisis se desprenden además tres conclusiones de diseño. Primero, el análisis de estabilidad responde la pregunta que planteaba el fabricante, que evaluaba omitir el larguero interior para facilitar la fabricación. La configuración sin larguero resulta claramente peor, por lo que el larguero es estructuralmente necesario; y dado que engrosarlo no mejora los multiplicadores, el cumplimiento del criterio exige rigidizar los paneles de piel del vano central. La comparación entre las dos vías disponibles es concluyente, ya que subdividir el panel más solicitado con un larguerillo longitudinal cuadruplica su capacidad por algo más de diez kilogramos, mientras que alcanzar el mismo resultado engrosando la piel exige del orden de cien. Segundo, las zonas que el diseño debe atender son distintas según el criterio, ya que la tensión máxima se localiza en las conexiones pala-brazo, lo que las señala como prioritarias para el detallamiento del laminado y las inspecciones en servicio, mientras que los modos críticos de pandeo se concentran en los paneles del vano central. Tercero, el diseño estructural de un rotor de eje vertical está gobernado por su velocidad máxima de giro, y este trabajo lo cuantifica, ya que la carga inercial permanente que impone la rotación produce por sí sola $114{,}5$~MPa, que son la mayor parte de la tensión del caso gobernante, frente a la fracción restante que aporta el viento en su condición extrema de diseño a 50 años. El dimensionamiento de la pala responde, por tanto, al valor que se fije para esa velocidad, $42$~rpm en este caso, equivalentes a $20{,}4\,g$ de aceleración centrífuga, y solo de forma secundaria al recurso eólico del emplazamiento. La consecuencia práctica es que fijar la velocidad máxima de giro es una decisión estructural antes que operacional, y que no puede resolverse mirando solo la curva de potencia.

En el plano metodológico, la comparación entre los DLC 1.1 y 1.3 deja una lección aplicable a cualquier verificación de este tipo. La Sección 7.4.2 de la norma IEC 61400-1~\parencite{IEC61400-1} permite omitir el análisis del DLC 1.1 cuando el DLC 1.3 lo envuelve, y ese envolvimiento no se produce aquí, ya que el DLC 1.3 gobierna en cuatro de las trece magnitudes evaluadas y en ninguna de las seis del apoyo superior, que es la sección más solicitada. Lo relevante es cómo se llega a esa conclusión. Si la comparación se hubiese resuelto sobre una única señal escalar, la fuerza total integrada sobre la pala, el DLC 1.3 la supera en un $5{,}4\%$ y el criterio habría autorizado prescindir precisamente del caso que gobierna. Evaluarlo magnitud a magnitud, sobre el conjunto de efectos de carga que la norma señala para la raíz de la pala, no es por tanto un refinamiento formal sino la diferencia entre acertar y equivocarse en la elección del caso de diseño. Que el ordenamiento obtenido coincida con el que \textcite{Galinos2015} reporta para otro rotor VAWT sugiere, además, que no se trata de una particularidad de esta máquina.

Una segunda lección metodológica proviene del sentido de la carga. El instante que maximiza el momento aerodinámico en la sección de apoyo resulta ser uno en que la fuerza normal apunta hacia el eje del rotor, de modo que se opone a la fuerza centrífuga en lugar de sumarse a ella. En una máquina donde la carga inercial domina la solicitación, el mayor extremo aerodinámico no produce el mayor estado tensional, y verificar solo ese instante habría dejado fuera la combinación realmente desfavorable. De ahí que cada componente del momento se extrapole en sus dos sentidos y que el modelo estructural reciba dos perfiles de carga por caso, uno por cada sentido de la componente normal. La advertencia es general para rotores de eje vertical, en los que la carga centrífuga es permanente y de signo fijo mientras la aerodinámica cambia de sentido a lo largo de cada revolución.

Una tercera observación atañe a la lectura del coeficiente de potencia. La definición convencional refiere la potencia media al cubo de la velocidad media del intervalo, supuesto válido con flujo uniforme y no con viento turbulento, donde la potencia disponible corresponde a la media del cubo de la velocidad. La diferencia no es menor a baja velocidad de viento, donde la intensidad de turbulencia relativa es mayor, y lleva a que el coeficiente así calculado alcance el límite de Betz sin que exista violación física alguna. Referido a la potencia que el campo transporta efectivamente, el máximo se sitúa en $0{,}503$ y se mantiene por debajo del límite en todo el rango. Cualquier comparación de rendimiento entre máquinas exige, por tanto, verificar antes sobre qué magnitud está definido el coeficiente que se compara.

El trabajo demuestra, además, que la metodología es replicable y de bajo costo de entrada, ya que la cadena completa, desde la caracterización del emplazamiento hasta la verificación por elementos finitos, se ejecutó con datos públicos chilenos y con QBlade, de código abierto, quedando Ansys como única herramienta comercial. Los contrastes independientes señalados más arriba son los que permiten sostener los resultados con la información disponible, pero no sustituyen a la medición. Si el proyecto avanza, la caracterización del emplazamiento debería respaldarse con una campaña de prospección eólica que mida in situ las condiciones de viento y turbulencia, ya que los parámetros que aquí se derivaron de reanálisis, entre ellos la velocidad media, la intensidad de turbulencia y la clase IEC del sitio, condicionan toda la cadena de cargas.

Esa campaña debería extenderse al estado de mar, que este trabajo no analizó por alcance. Su relevancia para las cargas de la pala no es la misma en todas las máquinas de esta familia. \textcite{Guo2018} simulan un H-rotor flotante bajo viento y oleaje combinados y encuentran que el movimiento de la plataforma amplifica las amplitudes de las fuerzas aerodinámicas hasta un $9{,}4\%$ en la componente normal y un $16{,}3\%$ en la tangencial respecto de la misma turbina sobre fundación fija, e incorpora a su contenido espectral la frecuencia del oleaje junto a la de paso de pala. \textcite{Borg2015} identifican la condición bajo la cual ese acoplamiento se vuelve crítico, ya que la excitación aerodinámica de una VAWT se concentra en su frecuencia de paso de pala, y en rotores muy grandes y de giro lento esa frecuencia cae dentro de la banda $0{,}5$--$0{,}9$~rad/s en que se concentra la energía del oleaje, con el consiguiente riesgo de excitación simultánea de origen aerodinámico e hidrodinámico.

Ese solapamiento no se produce en el rotor estudiado. Con tres palas girando a $42$~rpm, su frecuencia de paso de pala es de $13{,}2$~rad/s, más de un orden de magnitud por encima del extremo superior de esa banda, a diferencia del rotor de $5$~MW y giro lento que analizan \textcite{Borg2015}, cuya frecuencia de $0{,}77$~rad/s cae de lleno en ella. Cabe esperar, por tanto, que el oleaje incida sobre las cargas de la pala por la vía del movimiento de la plataforma, en el orden de magnitud que reporta \textcite{Guo2018}, y no por resonancia con la excitación del rotor. Comprobarlo exige el modelo acoplado que el alcance de este trabajo no incluyó.

\section{Limitaciones del alcance}

Los resultados deben interpretarse dentro del alcance definido al inicio del trabajo, que estableció el análisis bajo condiciones nominales de operación. Por ello se evaluaron los casos de producción de energía DLC 1.1, 1.2 y 1.3, quedando fuera del estudio las situaciones de arranque, parada, falla del sistema de control y condiciones de reposo, que la norma contempla en otros casos de carga. Esta delimitación es consistente con la evidencia de que en una VAWT la operación normal con turbulencia gobierna las cargas en la raíz de la pala~\parencite{Galinos2016,Sutherland2012}, pero no la convierte en una condición suficiente, ya que trabajos recientes sobre H-rotores rígidos y sobre plataformas flotantes han identificado contribuciones no despreciables en casos habitualmente descartados, en particular el rotor detenido bajo viento extremo~\parencite{Moore2024,Sakib2022}. Verificar la pala frente a ese conjunto completo excede el alcance comprometido y se plantea como extensión del estudio.

Dentro de ese marco, la verificación de la velocidad máxima de rotación adoptada ($42$~rpm) es estructural, esto es, acredita que la pala resiste las cargas centrífugas y aerodinámicas asociadas, pero no incluye el análisis modal del rotor ni la comprobación de resonancia entre la frecuencia de paso de pala y los modos propios de la estructura. Dicha verificación excede el dominio de cálculo definido, que comprende únicamente la pala, y requiere las propiedades de los brazos, la torre y la plataforma. La experiencia documentada muestra que ese criterio no es accesorio, ya que en las máquinas construidas es la estructura, y no la aerodinámica, la que termina fijando la velocidad máxima de giro. En el prototipo de Sandia, el cruce de la frecuencia de paso de pala con el primer modo de torre en el plano acotó la velocidad máxima practicable a $38$~rpm~\parencite{Sutherland2012}, y en la T1-turbine el límite operacional se redujo de $33$ a $22$~rpm tras la instalación de los tensores de la torre, también por resonancia~\parencite{Mollerstrom2017}. Su comprobación constituye, por tanto, un requisito previo a la adopción definitiva de las $42$~rpm.

Conviene explicitar, en relación con lo anterior, que la velocidad máxima de rotación se trató aquí como un dato de entrada y no como una variable de diseño. El valor se fijó a partir del criterio aerodinámico, aplicando el TSR de diseño a la velocidad de viento nominal, y sobre él se verificó la pala en un único sentido, sin iterar entre la velocidad adoptada y el dimensionamiento resultante. Esa decisión responde al alcance comprometido, que es verificar una configuración dada y no optimizarla, pero sobre todo a la información disponible, ya que el espesor de pared no proviene del fabricante sino que se estimó a partir de la masa total de la pala. Iterar espesores sobre una magnitud que es en sí misma un supuesto habría entregado un óptimo aparente, no un diseño mejor. Los resultados respaldan además que la velocidad de giro es el parámetro estructuralmente decisivo, ya que la carga inercial que impone aporta el $70\%$ de la tensión del caso gobernante, de modo que una eventual iteración entre velocidad de rotación y dimensionamiento de la sección es una línea pertinente, condicionada a disponer de la definición del laminado y del modelo modal del rotor completo.

La influencia del oleaje y de la dinámica acoplada entre el rotor y la plataforma tampoco se consideró, ya que su modelación requiere información de la plataforma flotante y de los sistemas de fondeo con la que no se cuenta. El desarrollo de este trabajo coincidió con la fase de aprobación de presupuesto del proyecto, por lo que parte de la información de diseño aún no estaba disponible al término de la práctica; las condiciones no especificadas, entre ellas la velocidad máxima de rotación, el exponente de fatiga del laminado y el espesor de pared uniforme, se asumieron a partir de los datos públicos del fabricante y de valores documentados en la literatura, dejando cada supuesto declarado explícitamente en la metodología. Que los resultados aerodinámicos y estructurales resultaran consistentes con la evidencia experimental y con turbinas de escala comparable respalda la razonabilidad de estos supuestos.

Por último, el modelo estructural representa el laminado mediante el conjunto isótropo equivalente declarado en la Sección~\ref{sec:material_pala}, idealización cuyo alcance ya se discutió al comentar el resultado de estabilidad y que afecta sobre todo a esa verificación.

\section{Trabajo futuro}

De las limitaciones anteriores se desprenden extensiones naturales del estudio. Cuando el proyecto disponga de los parámetros reales de funcionamiento de la turbina, esto es, la velocidad máxima de rotación de diseño, el laminado definitivo y la distribución de espesores, el análisis puede repetirse refinando los supuestos, con la ventaja de que la cadena de cálculo ya está construida y validada.

Una extensión de menor coste y de efecto inmediato sobre la confianza en los resultados es ampliar el número de realizaciones por velocidad en la banda alta. La Sección G.4.1 del Anexo G pide al menos quince realizaciones por velocidad entre $V_r - 2$~m/s y la velocidad de corte, frente a las seis con que se ejecutó esta campaña~\parencite{IEC61400-1}. Son precisamente las velocidades que dominan la cola de la extrapolación y la práctica totalidad del daño por fatiga, de modo que el refuerzo del muestreo ahí reduce la incertidumbre donde más pesa, con un costo de cómputo acotado a siete velocidades.

Un refinamiento adicional del modelo estructural consiste en simular el laminado capa por capa mediante el módulo ACP de Ansys, en lugar de evaluar el compuesto como material homogéneo, lo que permitiría identificar el modo de falla de la fibra y capturar la ortotropía que condiciona tanto la resistencia como la carga crítica de pandeo. Con un margen en tensión del $14\%$ y una verificación de estabilidad no satisfecha, ese nivel de detalle gana relevancia. La herramienta es conocida y su aplicación quedó supeditada únicamente a disponer de la definición del laminado.

En el plano de la estabilidad, el incumplimiento de la verificación de pandeo define la extensión prioritaria y orienta el rediseño de la sección. El margen disponible en fatiga indica que el material no está limitado por la acumulación de daño, de modo que todo refuerzo adicional debe destinarse a rigidizar los paneles de piel del vano central, cuya capacidad gobierna el modo crítico una vez incorporado el larguero: un segundo larguero, larguerillos longitudinales, costillas transversales o un mayor espesor local. Sobre la configuración corregida corresponde repetir la cadena de verificación completa, dado que cualquier rigidización modifica la masa, la rigidez y la distribución de tensiones, y complementarla con un análisis de pandeo no lineal que incorpore imperfecciones geométricas, pues el método lineal por autovalores entrega una cota superior de la carga crítica.

Junto a lo anterior, el análisis modal del rotor completo permitiría verificar que la velocidad de rotación adoptada no induzca resonancia, condición que en la turbina de referencia obligó a reducir el límite operacional. A ello se suma someter la pala a los restantes casos de carga de IEC 61400-1, en particular los eventos transitorios de arranque y parada y las condiciones extremas con rotor detenido, y la incorporación del oleaje y de la dinámica acoplada entre el rotor y la plataforma, una vez disponible la información del sistema de flotación, que permitiría verificar la pala en el entorno offshore completo.

Todas estas extensiones comparten una condición favorable: la cadena de análisis quedó construida, contrastada en cada etapa y documentada, de modo que incorporarlas constituye un refinamiento y no una reconstrucción. El trabajo cierra con dos resultados de distinta naturaleza. Por una parte, un procedimiento de verificación estructural completo, que abarca desde la caracterización del emplazamiento hasta el modelo de elementos finitos y resulta aplicable a otras palas de configuración similar. Por otra, un diagnóstico de la pala analizada, en el que la resistencia última, la fatiga y las deflexiones quedan verificadas, mientras que la estabilidad de los paneles de piel constituye la condición pendiente para considerar el diseño apto para fabricación.
